\documentclass[12pt,a4paper]{article}
\usepackage[utf8]{inputenc}
\usepackage[T1]{fontenc}
\usepackage[spanish]{babel}
\usepackage{lmodern}
\usepackage{geometry}
\usepackage{hyperref}
\usepackage{enumitem}
\usepackage{parskip}

\geometry{margin=2.5cm}
\hypersetup{
  colorlinks=true,
  linkcolor=black,
  urlcolor=blue,
  pdftitle={Informe del proyecto The Daily Scribble},
  pdfauthor={Harry \& EPN}
}

\title{Informe del Proyecto\\\large The Daily Scribble}
\author{Harry \& EPN}
\date{30 de marzo de 2026}

\begin{document}

\maketitle

\section{Objetivo del proyecto}
El objetivo de \textbf{The Daily Scribble} es ofrecer un espacio creativo y ligero para dibujar en sesiones cortas, sin presion y con enfoque diario. La aplicacion busca incentivar la practica constante mediante una \textit{palabra del dia}, un tiempo limite configurable y una experiencia de publicacion simple en una galeria local.

\section{Que hace la aplicacion}
La aplicacion permite que cualquier persona:
\begin{itemize}[leftmargin=*]
  \item Elija una sesion de dibujo de 1, 5, 10 o 15 minutos.
  \item Dibuje sobre un lienzo usando multiples herramientas.
  \item Firme su obra de manera opcional al momento de compartir.
  \item Descargue su dibujo en formato PNG.
  \item Visualice dibujos compartidos en una galeria con animacion fisica.
\end{itemize}

Ademas, el contenido es \textbf{efimero por dia}: el estado se asocia a la fecha actual y se reinicia cuando cambia el dia.

\section{Funcionalidades principales}
\begin{itemize}[leftmargin=*]
  \item \textbf{Palabra diaria automatica}: se calcula en funcion de la fecha.
  \item \textbf{Temporizador de sesion}: controla el tiempo restante y bloquea la edicion al finalizar.
  \item \textbf{Herramientas de dibujo}: pincel, borrador, linea, rectangulo, circulo, bote de pintura y gotero.
  \item \textbf{Control de edicion}: deshacer y rehacer por historial de estados del canvas.
  \item \textbf{Gestion de color}: paleta base, selector personalizado y agregado de colores.
  \item \textbf{Exportacion}: descarga de imagen en PNG con firma opcional.
  \item \textbf{Galeria local}: visualizacion de resultados por duracion con motor fisico Matter.js.
\end{itemize}

\section{Como esta hecho (arquitectura tecnica)}
El proyecto es una aplicacion web \textbf{estatica del lado del cliente} (sin backend), organizada en tres archivos:

\begin{itemize}[leftmargin=*]
  \item \textbf{index.html}: define las vistas (inicio, galeria, preparacion, canvas) y los modales.
  \item \textbf{style.css}: contiene el diseno visual, estilos responsivos y componentes de interfaz.
  \item \textbf{app.js}: implementa toda la logica de negocio, dibujo, temporizador, persistencia y galeria.
\end{itemize}

\subsection{Flujo funcional}
\begin{enumerate}[leftmargin=*]
  \item El usuario entra a la vista de inicio.
  \item Selecciona \textit{Empezar a Dibujar} y escoge una duracion.
  \item Se activa la sesion de dibujo con temporizador.
  \item Al terminar, puede firmar y compartir.
  \item La obra se guarda en estado local y aparece en la galeria correspondiente.
\end{enumerate}

\subsection{Persistencia y estado}
La aplicacion usa \texttt{localStorage} con una clave diaria para guardar:
\begin{itemize}[leftmargin=*]
  \item Fecha actual.
  \item Palabra diaria.
  \item Dibujos de galeria clasificados por duracion.
\end{itemize}
Esto permite mantener la experiencia de "obra del dia" sin servidor externo.

\section{Tecnologias utilizadas}
\begin{itemize}[leftmargin=*]
  \item HTML5
  \item CSS3
  \item JavaScript (Vanilla)
  \item Canvas API
  \item Matter.js (cargado desde CDN)
\end{itemize}

\section{Ventajas del enfoque actual}
\begin{itemize}[leftmargin=*]
  \item Despliegue sencillo (ideal para GitHub Pages).
  \item Sin dependencias de backend ni base de datos.
  \item Rendimiento rapido en navegadores modernos.
  \item Codigo facil de entender y mantener para un MVP.
\end{itemize}

\section{Posibles mejoras futuras}
\begin{itemize}[leftmargin=*]
  \item Galeria compartida real con backend y almacenamiento en la nube.
  \item Sistema de cuentas o identidad persistente.
  \item Compatibilidad PWA para uso offline.
  \item Herramientas avanzadas (capas, texto, relleno inteligente).
  \item Moderacion y filtros de contenido para galeria publica.
\end{itemize}

\section{Repositorio}
Codigo fuente disponible en:
\href{https://github.com/Harry-GNS/Scrible}{https://github.com/Harry-GNS/Scrible}

\end{document}
